\Needspace{8\baselineskip}
\Article{188}
Since all forms in the same period are properly equivalent, the question arises whether forms in different periods can also be properly equivalent. But before showing that \textit{this is impossible}, we must discuss some points about transformations of reduced forms.

Since we will often be dealing with transformations of forms in what follows, from now on we will always use the following abbreviated notation to keep the writing as short as possible. If a form $LXX+2MXY+NYY$ is transformed by the substitution $X=\alpha x+\beta y$, $Y=\gamma x+\delta y$ into the form $lxx+2mxy+nyy$, we will simply say that $(L,M,N)$ is transformed into $(l,m,n)$ by the substitution $\alpha,\beta,\gamma,\delta$. This way we will not need separate symbols for the indeterminates of each form under consideration. --- But clearly the \textit{first} indeterminate must be carefully distinguished from the \textit{second} in every form.

{\rightskip=0pt plus 4em\relax
Let the reduced form $(a,b,-a')\mathinner{\ldotp\ldotp\ldotp} f$ with determinant $D$ be given. As in Article~186, form the sequence of reduced forms infinite in both directions, $\mathinner{\ldotp\ldotp\ldotp\ldotp} {}''f,{}'f,f,f',f''\mathinner{\ldotp\ldotp\ldotp\ldotp}$, and let
\SourceLeaf{177}
\[
\begin{aligned}
f'&=(-a',b',a''),& f''&=(a'',b'',-a''')\quad\text{etc.}\\
{}'f&=(-{}'a,{}'b,a),& {}''f&=({}''a,{}''b,-{}'a)\quad\text{etc.}
\end{aligned}
\]
Set
\[
\begin{array}{llll}
\dfrac{b+b'}{-a'}=h',&\dfrac{b'+b''}{a''}=h'',&\dfrac{b''+b'''}{-a'''}=h'''&\text{etc.}\\[6pt]
\dfrac{{}'b+b}{a}=h,&\dfrac{{}''b+{}'b}{-{}'a}={}'h,&\dfrac{{}'''b+{}''b}{{}''a}={}''h&\text{etc.}
\end{array}
\]
Then clearly, if (as in Article~177) the numbers $\alpha',\alpha'',\alpha'''$ etc., $\beta',\beta'',\beta'''$ etc., and so on are formed by the following algorithm
{\small
\setlength{\arraycolsep}{3pt}
\[
\begin{array}{llll}
\alpha'=0&\beta'=-1&\gamma'=1&\delta'=h'\\
\alpha''=\beta'&\beta''=h''\beta'&\gamma''=\delta&\delta''=h''\delta'-1\\
\alpha'''=\beta''&\beta'''=h'''\beta''-\beta'&\gamma'''=\delta''&\delta'''=h'''\delta''-\delta'\\
\alpha''''=\beta'''&\beta''''=h''''\beta'''-\beta''&\gamma''''=\delta'''&\delta''''=h''''\delta'''-\delta''
\\
\multicolumn{4}{c}{\text{etc.}}
\end{array}
\]
}
$f$ will be transformed
\[
\begin{array}{@{}rlll@{}}
\text{into}&f'&\text{by the substitution}&\alpha',\ \beta',\ \gamma',\ \delta'\\
&f''&\mathinner{\ldotp\ldotp\ldotp\ldotp\ldotp\ldotp\ldotp\ldotp}&\alpha'',\ \beta'',\ \gamma'',\ \delta''\\
&f'''&\mathinner{\ldotp\ldotp\ldotp\ldotp\ldotp\ldotp\ldotp\ldotp}&\alpha''',\ \beta''',\ \gamma''',\ \delta'''
\end{array}
\]
\centerline{etc.}
and all these transformations will be proper.\par}

\Needspace{24\baselineskip}
Since ${}'f$ goes into $f$ by the proper substitution $0,-1,1,h$ (Article~158), $f$ goes into ${}'f$ by the proper substitution $h,1,-1,0$. For a similar reason ${}'f$ goes into ${}''f$ by the proper substitution ${}'h,1,-1,0$; ${}''f$ into ${}'''f$ by the proper substitution ${}''h,1,-1,0$, and so on. From Article~159, in the same way as in Article~177, we conclude that if the numbers ${}'\alpha,{}''\alpha,{}'''\alpha$ etc., ${}'\beta,{}''\beta,{}'''\beta$ etc., and so on are formed by the following algorithm
{\small
\setlength{\arraycolsep}{3pt}
\[
\begin{array}{llll}
{}'\alpha=h&{}'\beta=1&{}'\gamma=-1&{}'\delta=0\\
{}''\alpha={}'h\,{}'\alpha-1&{}''\beta={}'\alpha&{}''\gamma={}'h\,{}'\gamma&{}''\delta={}'\gamma\\
{}'''\alpha={}''h\,{}''\alpha-{}'\alpha&{}'''\beta={}''\alpha&{}'''\gamma={}''h\,{}''\gamma-{}'\gamma&{}'''\delta={}''\gamma\\
{}''''\alpha={}'''h\,{}'''\alpha-{}''\alpha&{}''''\beta={}'''\alpha&{}''''\gamma={}'''h\,{}'''\gamma-{}''\gamma&{}''''\delta={}'''\gamma
\\
\multicolumn{4}{c}{\text{etc.}}
\end{array}
\]
}
$f$ will be transformed
\SourceLeaf{178}
\[
\begin{array}{@{}rlll@{}}
\text{into}&{}'f&\text{by the substitution}&{}'\alpha,\ {}'\beta,\ {}'\gamma,\ {}'\delta\\
&{}''f&\mathinner{\ldotp\ldotp\ldotp\ldotp\ldotp\ldotp\ldotp\ldotp}&{}''\alpha,\ {}''\beta,\ {}''\gamma,\ {}''\delta\\
&{}'''f&\mathinner{\ldotp\ldotp\ldotp\ldotp\ldotp\ldotp\ldotp\ldotp}&{}'''\alpha,\ {}'''\beta,\ {}'''\gamma,\ {}'''\delta
\end{array}
\]
\centerline{etc.}
and all these transformations will be proper.

If we set $\alpha=1$, $\beta=0$, $\gamma=0$, $\delta=1$, these numbers will have the same relation to the form $f$ as $\alpha',\beta',\gamma',\delta'$ have to $f'$, $\alpha'',\beta'',\gamma'',\delta''$ to $f''$ etc., ${}'\alpha,{}'\beta,{}'\gamma,{}'\delta$ to ${}'f$ etc. Namely, the substitution $\alpha,\beta,\gamma,\delta$ takes the form $f$ into $f$. The infinite sequences $\alpha',\alpha'',\alpha'''$ etc., ${}'\alpha,{}''\alpha,{}'''\alpha$ etc. join together by inserting the term $\alpha$ so that they can be regarded as one continuous sequence, infinite in both directions and following the same law throughout: $\mathinner{\ldotp\ldotp\ldotp} {}'''\alpha,{}''\alpha,{}'\alpha,\alpha,\alpha',\alpha'',\alpha'''\mathinner{\ldotp\ldotp\ldotp}$. The law of this sequence is
\[
\begin{gathered}
{}'''\alpha+{}'\alpha={}''h\,{}''\alpha,\quad {}''\alpha+\alpha={}'h\,{}'\alpha,\quad {}'\alpha+\alpha'=h\alpha,\\
\alpha+\alpha''=h'\alpha',\quad \alpha'+\alpha'''=h''\alpha''\text{ etc.}
\end{gathered}
\]
or, in general (if we take a negative index written on the right to mean the same as a positive one on the left),
\[
\alpha^{m-1}+\alpha^{m+1}=h^m{\alpha^m}
\]
Similarly the sequence $\mathinner{\ldotp\ldotp\ldotp} {}''\beta,{}'\beta,\beta,\beta',\beta''\mathinner{\ldotp\ldotp\ldotp}$ is continuous, with the law
\[
\beta^{m-1}+\beta^{m+1}=h^{m+1}\beta^m
\]
and is essentially identical to the preceding sequence, with all its terms moved one place: ${}''\beta={}'\alpha$, ${}'\beta=\alpha$, $\beta=\alpha'$ etc. The law of the continuous sequence $\mathinner{\ldotp\ldotp\ldotp\ldotp} {}''\gamma,\allowbreak {}'\gamma,\allowbreak \gamma,\allowbreak \gamma',\allowbreak \gamma''\mathinner{\ldotp\ldotp\ldotp\ldotp\ldotp}$ is
\[
\gamma^{m-1}+\gamma^{m+1}=h^m{\gamma^m}
\]
and the law of $\mathinner{\ldotp\ldotp\ldotp} {}''\delta,{}'\delta,\delta,\delta',\delta''\mathinner{\ldotp\ldotp\ldotp}$ is
\[
\delta^{m-1}+\delta^{m+1}=h^{m+1}\delta^m
\]
and in addition, in general, $\delta^m=\gamma^{m+1}$.

\textit{Example.} Let the given form $f$ be $(3,8,-5)$. It will be transformed
\SourceLeaf{179}
\[
\setlength{\arraycolsep}{4pt}
\begin{array}{@{}lll@{}}
\multicolumn{2}{l}{\text{into the form}}&\text{by the substitution}\\
{}'''''''f&(-10,7,3)&-805,\ -152,\ +143,\ +27\\
{}''''''f&(3,8,-5)&-152,\ +45,\ +27,\ -8\\
{}'''''f&(-5,7,6)&+45,\ +17,\ -8,\ -3\\
{}''''f&(6,5,-9)&+17,\ -11,\ -3,\ +2\\
{}'''f&(-9,4,7)&-11,\ -6,\ +2,\ +1\\
{}''f&(7,3,-10)&-6,\ +5,\ +1,\ -1\\
{}'f&(-10,7,3)&+5,\ +1,\ -1,\ 0\\
f&(3,8,-5)&+1,\ 0,\ 0,\ +1\\
f'&(-5,7,6)&0,\ -1,\ +1,\ -3\\
f''&(6,5,-9)&-1,\ -2,\ -3,\ -7\\
f'''&(-9,4,7)&-2,\ +3,\ -7,\ +10\\
f''''&(7,3,-10)&+3,\ +5,\ +10,\ +17\\
f'''''&(-10,7,3)&+5,\ -8,\ +17,\ -27\\
f''''''&(3,8,-5)&-8,\ -45,\ -27,\ -152\\
f'''''''&(-5,7,6)&-45,\ +143,\ -152,\ +483
\end{array}
\]
\centerline{etc.}

\Needspace{8\baselineskip}
\Article{189}

The following points about this algorithm should be noted.

1) All $a,a',a''$ etc., ${}'a,{}''a$ etc. have the same sign; all $b,b',b''$ etc., ${}'b,{}''b$ etc. are positive. In the sequence $\mathinner{\ldotp\ldotp\ldotp} {}''h,{}'h,h,h',h''\mathinner{\ldotp\ldotp\ldotp\ldotp}$ the signs alternate: if all $a,a'$ etc. are positive, $h^m$ or ${}^m h$ is positive when $m$ is even and negative when $m$ is odd; but if $a,a'$ etc. are negative, $h^m$ or ${}^m h$ is negative for even $m$ and positive for odd.

2) If $a$ is positive, and so $h'$ negative, $h''$ positive, etc., then $\alpha''=-1$ is negative, $\alpha'''=h''\alpha''$ negative and $>\alpha''$ (or $=\alpha''$ if $h''=1$); $\alpha''''=h'''\alpha'''-\alpha''$ positive and $>\alpha'''$ (since $h'''\alpha'''$ is positive and $\alpha''$ negative); $\alpha'''''=h''''\alpha''''-\alpha'''$ positive and $>\alpha''''$ (since $h''''\alpha''''$ is positive), and so on. We easily conclude that the sequence $\alpha',\alpha'',\alpha'''$ etc. grows without bound, and that two positive signs are always followed by two negative ones, so that $\alpha^m$ has the sign $+,+,-,-$ according as $m\equiv0,1,2,3\ (\mathrm{mod.}\ 4)$. If $a$ is negative, a similar argument gives $\alpha''$ negative, $\alpha'''$ positive and either $>$ or $=\alpha''$; $\alpha''''$ positive and $>\alpha'''$; $\alpha'''''$ negative and $>\alpha''''$, and so on,
\SourceLeaf{180}
so that the sequence $\alpha',\alpha'',\alpha'''$ etc. keeps growing, and the term $\alpha^m$ has the sign $+,-,-,+$ according as $m\equiv0,1,2,3\ (\mathrm{mod.}\ 4)$.

\Needspace{32\baselineskip}
3) In this way we find that all four infinite sequences $\alpha',\alpha'',\alpha'''$ etc., $\gamma,\gamma',\gamma''$ etc.; $\alpha',\alpha,{}'\alpha,{}''\alpha$ etc.; $\gamma,{}'\gamma,{}''\gamma$ etc. keep growing, and so do the following sequences identical to them: $\beta,\beta',\beta''$ etc.; ${}'\delta,\delta,\delta',\delta''$ etc.; $\beta,{}'\beta,{}''\beta$ etc.; ${}'\delta,{}''\delta$ etc. According as $m\equiv0,1,2,3\ (\mathrm{mod.}\ 4)$, the signs are
\[
\begin{array}{ll}
\text{of }\alpha^m,\quad +\ \pm\ -\ \mp;&\text{of }\beta^m,\quad \pm\ -\ \mp\ +\\[6pt]
\text{of }\gamma^m,\quad \pm\ +\ \mp\ -;&\text{of }\delta^m,\quad +\ \mp\ -\ \pm\\[6pt]
\text{of }{}^m\alpha,\quad +\ \pm\ -\ \mp;&\text{of }{}^m\beta,\quad \mp\ +\ \pm\ -\\[6pt]
\text{of }{}^m\gamma,\quad \mp\ -\ \pm\ +;&\text{of }{}^m\delta,\quad +\ \mp\ -\ \pm
\end{array}
\]
with the upper signs applying when $a$ is positive and the lower when $a$ is negative. Keep the following property especially in mind: if $m$ denotes any positive index, $\alpha^m$ and $\gamma^m$ have the same sign when $a$ is positive. Opposite signs when $a$ is negative, and similarly for $\beta^m$ and $\delta^m$; conversely, ${}^m\alpha$ and ${}^m\gamma$, or ${}^m\beta$ and ${}^m\delta$, have the same sign when $a$ is negative and opposite signs when $a$ is positive.

4) In the notation of Article~27, the magnitude of $\alpha^m$ etc. can be expressed conveniently as follows. Set
\[
\begin{gathered}
\mp h'=k',\quad \pm h''=k'',\quad \mp h'''=k'''\text{ etc.}\\
\pm h=k,\quad \mp{}'h={}'k,\quad \pm{}''h={}''k\text{ etc.}
\end{gathered}
\]
so that all $k',k''$ etc., $k,{}'k$ etc. are positive numbers:
{\small
\setlength{\arraycolsep}{3pt}
\[
\begin{array}{ll}
\alpha^m=\pm[k'',k''',k''''\mathinner{\ldotp\ldotp\ldotp\ldotp} k^{m-1}];&\beta^m=\pm[k'',k''',k''''\mathinner{\ldotp\ldotp\ldotp} {k^m]}\\[6pt]
\gamma^m=\pm[k',k'',k'''\mathinner{\ldotp\ldotp\ldotp\ldotp\ldotp} k^{m-1}];&\delta^m=\pm[k',k'',k'''\mathinner{\ldotp\ldotp\ldotp\ldotp} k^m]\\[6pt]
{}^m\alpha=\pm[k,{}'k,{}''k\mathinner{\ldotp\ldotp\ldotp\ldotp\ldotp\ldotp}{}^{m-1}k];&{}^m\beta=\pm[k,{}'k,{}''k\mathinner{\ldotp\ldotp\ldotp\ldotp}{}^{m-2}k]\\[6pt]
{}^m\gamma=\pm[{}'k,{}''k\mathinner{\ldotp\ldotp\ldotp\ldotp\ldotp\ldotp}{}^{m-1}k];&{}^m\delta=\pm[{}'k,{}''k\mathinner{\ldotp\ldotp\ldotp\ldotp}{}^{m-2}k]
\end{array}
\]
}
The \textit{signs}, however, must be determined by the rules just given. With these formulas, whose proof we omit because it is easy, the calculation can always be carried out very quickly.
\SourceLeaf{181}

\Needspace{8\baselineskip}
\Article{190}

\textsc{Lemma.} \textit{Let $m,\mu,m',n,\nu,n'$ denote any integers, but with none of the last three $=0$. I say that if $\frac{\mu}{\nu}$ lies strictly between $\frac{m}{n}$ and $\frac{m'}{n'}$, and $mn'-nm'={\mp1}$, the denominator $\nu$ is greater than $n$ and $n'$.}

\textit{Proof.} Clearly $\mu nn'$ lies between $\nu mn'$ and $\nu nm'$, and so differs from either bound by less than the bounds differ from each other; \textit{that is}, $\nu mn'-\nu nm'>\mu nn'-\nu mn'$ and $>\mu nn'-\nu nm'$, or $\nu>n'(\mu n-\nu m)$ and $>n(\mu n'-\nu m')$. Since $\mu n-\nu m$ is certainly not $=0$ (otherwise $\frac{\mu}{\nu}=\frac{m}{n}$, contrary to the hypothesis), nor is $\mu n'-\nu m'=0$ (for a similar reason), but each is at least $=1$, it follows that $\nu>n'$ and $>n$. \mbox{\textit{Q. E. D.}}

So $\nu$ clearly cannot be $=1$; \textit{that is}, if $mn'-nm'=\pm1$, no integer can lie between the fractions $\frac{m}{n}$, $\frac{m'}{n'}$. In particular, zero cannot lie between them, \textit{that is}, these fractions cannot have opposite signs.
